\section{结论}

我们围绕``算力约束下提升大语言模型能力的资源配置''建立了完整决策链：
问题一给出可解释、可迁移的数据质量评分与配比效应模型，问题二把质量维度
嵌入经典标度律得到广义能力方程，问题三在算力约束下求解最优资源配置并
给出结构性转移的数学刻画，问题四把``能力--资源''关系用于行业前沿分解
与预测。四问结论汇总于表 \ref{tab:conclusion}。

\begin{table}[htbp]
\centering
\caption{四问结论总览}
\label{tab:conclusion}
\small
\begin{tabular}{>{\raggedright\arraybackslash}p{1.7cm}|>{\raggedright\arraybackslash}p{3.6cm}|>{\raggedright\arraybackslash}p{4.7cm}}
\toprule
问题 & 最优方案 & 关键数字 \\
\midrule
一 & 熵-CRITIC 组合赋权 + TOPSIS；对称截尾消解；17 域岭回归 & 域级质量 book 0.631 居首；配比收缩 $\kappa=0.314031$；调节收益约 3\% \\
二 & 质量--参数量交互广义标度律 & $R^2=0.979073$；$\varepsilon_Q=0.146080$；质量 0.1 等价 0.288515B 参数 \\
三 & 三通道成本 + 多框架 SLSQP 联合优化 & 收益 0.8\%--6.7\%；$N^{*}\sim C^{0.46}$；转移阈值 $10^{18}$ 量级 \\
四 & 90\% 分位数回归 + 双情景预测 & 规模贡献 83.7\%；$S_{12}=71.7$ / $S_{24}=123.9$（放缓 57.1/78.6） \\
\bottomrule
\end{tabular}
\end{table}

方法论上有三点体会。其一，结论的稳健性靠交叉验证而非断言：四个问题均用
多种独立框架/方法复算（评分八方法、标度律六拟合框架、优化九求解路径、
前沿五框架），全部收敛到一致结果；其二，数据口径要放在明处：我们诊断出
附件 B8 的 $Q$ 方向与 B6/B7 相反、C3 前沿含孤立高值等问题并给出规避与
对照处理，任何结论都标注了可信度边界；其三，中间过程应当可复核：拟合、
优化与预测的全部分布数据、Bootstrap 区间与结果文件均可逐步核对，
结论因此可查、可复现。

最后说明适用边界。全部结论建立在赛题附件数据之上，广义标度律的质量项
参数受 B 族半合成实验口径限制，$\ln S\sim\ln N+t$ 的线性外推在较长预测
跨度后可能因算法饱和失效；若引入真实多质量档预训练数据或滚动再估计
窗口，模型可平滑升级，既有框架无需改动。

\subsection{与题面要求的逐项对照}

为便于评审核对，把四问的题面要求与本团队的交付逐项对照（表
\ref{tab:checklist}），覆盖质量评分、广义标度律、最优配置与前沿预测四类
输出，全部交付均含对应表格、公式或附录定位。

\begin{table}[htbp]
\centering
\caption{题面要求与本团队交付的逐项对照}
\label{tab:checklist}
\small
\begin{tabular}{>{\raggedright\arraybackslash}p{1.7cm}|>{\raggedright\arraybackslash}p{5.2cm}|>{\raggedright\arraybackslash}p{3.4cm}}
\toprule
问题 & 题面要求 & 交付 \\
\midrule
一 & 质量评分方法与域级质量；冲突消解规则；配比--损失模型及规模转移 & 样本级/域级质量分（表 \ref{tab:p1_domain}）；双族冲突指数与对称截尾（第 5 节）；17 域岭回归与 $\kappa$ 收缩律 \\
二 & 含质量维度广义标度律；参数估计；跨家族泛化验证；弹性与替代关系 & 交互形式（式 \ref{eq:gen_choice}，$R^2=0.979073$）；家族偏移修正验证 $R^2=0.830$；弹性排序与等价参数（表 \ref{tab:p2_forms}） \\
三 & 三档预算最优配置；结构性转移的数学定义与识别；灵敏度 & SLSQP 最优解与 $Q^{*}\to1$；成本份额/边界双指标 + KKT 识别（第 7 节）；阈值口径敏感性 \\
四 & 规模/非规模贡献分解；12/24 个月前沿预测含不确定性；C8 逐任务聚合 & 规模贡献 83.7\%；双情景 90\% 区间（表 \ref{tab:backtest}）；MATH/IFEval 逐任务增速 \\
全文 & 结果可复现；口径一致 & 脚本随附、六位小数口径、附录 B 数据口径说明 \\
\bottomrule
\end{tabular}
\end{table}
